% =====================================================================
% 11_faddeev_solver.tex
% 第 11 章 Faddeev 求解器：Neumann 级数 + Padé 重求和
% Wave 3 / Agent C
% =====================================================================

% 章节内局部宏（仅本章使用，遵守"不修改 physics_macros"约束）
\providecommand{\Aker}{\mathbf{A}}                    % 驱动核 A = (C^T) P V C
\providecommand{\Kker}{\mathbf{K}}                    % 一阶 Neumann 核 K = G A
\providecommand{\Gker}{\mathbf{G}}                    % 通道分辨算符的对角矩阵
\providecommand{\Iker}{\mathbf{I}}                    % 单位阵
\providecommand{\Cker}{\mathbf{C}}                    % FWP→SWP 基变换
\providecommand{\Pker}{\mathbf{P}}                    % 离散 P123
\providecommand{\Vker}{\mathbf{V}}                    % 两体势矩阵
\providecommand{\Uker}{\mathbf{U}}                    % U-振幅块
\providecommand{\Lker}{\mathbf{L}}                    % LHS 矩阵 (1 - GA)
\providecommand{\Rker}{\mathbf{R}}                    % RHS 矩阵 A
\providecommand{\Pade}[2]{P_{[#1/#2]}}                % Padé 表 [N/M]
\providecommand{\rConv}{R_{\text{conv}}}              % 收敛半径
\providecommand{\PVon}{(\Pop\Vop)}                    % 论文里的 PV 复合算符

\chapter{Faddeev 求解器：Neumann 级数 + Padé 重求和}
\label{ch:11_faddeev_solver}

% =========================================
\section{物理动机}
\label{sec:ch11-motivation}

\paragraph{从 AGS 方程到一台数值机器。}
\cref{ch:04_faddeev_ags} 把三体散射归结为 AGS 方程
\begin{equation}
\label{eq:ch11-ags-recap}
\Uop \;=\; \PVon\,(\Iker - \Gzero\,\PVon)^{-1},
\end{equation}
\cref{ch:05_wpcd_math,ch:06_wp_swp_states,ch:07_pw_symm_states,ch:08_potential_matrix,ch:09_permutation_p123,ch:10_resolvent}
则把右端的每一个对象——\(\Vop, \Pop, \Gzero\) 以及 SWP 基矢的 \(\Cker\)
变换矩阵——拆成了\textbf{有限维矩阵}。本章要做的事，就是从这些
矩阵装配出一个\emph{具体可运行}的求解器：\emph{即给定全部矩阵元，
我们如何把 \cref{eq:ch11-ags-recap} 的逆矩阵作用到驱动项上，得到
on-shell 的 \(\Uampl{\PWchanp}{\PWchan}\)？}

把上式改写成更接近代码的形式。引入\textbf{驱动核}\(\Aker = \Cker^\top \Pker \Vker \Cker\)
（在 SWP 基下）和\textbf{一阶 Neumann 核}\(\Kker = \Gker\,\Aker\)，方程变为
\begin{equation}
\label{eq:ch11-fixed-point}
\Uker \;=\; \Aker + \Kker\,\Uker
\;=\; (\Iker - \Kker)^{-1}\,\Aker.
\end{equation}
\(\Uker\) 是一个稠密块——但物理上我们\emph{只需要}它的若干 on-shell 行
（\(\PWchanp\)：deuteron-out, \(q\)：质心动量），相邻的近 on-shell 与 off-shell
元素只是中间数据。这一观察直接决定了下面所有算法选择。

\paragraph{两条求解路径：dense vs.\ Neumann+Padé。}
\cref{eq:ch11-fixed-point} 是个线性方程组，原则上有两种求解办法：
\begin{itemize}
  \item \textbf{Dense 路径}（\texttt{solve\_dense=true}）：直接构造满矩阵
        \((\Iker - \Kker)\)，调用 LAPACK 的 \texttt{LU+TRS} 一次性求逆。
        优点：数值上等价于"把方程严格解出来"，没有迭代误差，可作为
        \emph{benchmark oracle}；缺点：内存与时间均为 \(\mathcal{O}(D^3)\)，
        其中 \(D = \Nalpha \cdot \NqWP \cdot \NpWP\) 通常达到 \(10^4\sim 10^5\)，
        导致 \(D^2\) 浮点数（\(\sim\)千 GiB）根本进不了内存——
        因此只能跑 \emph{toy} 规模（\(\Np=\Nq=10\sim 12\)）。
  \item \textbf{Neumann + Padé 路径}（默认）：把
        \((\Iker - \Kker)^{-1}\) 展成 Neumann 级数
        \(\sum_{n=0}^{\infty}\Kker^n\)，每生成一项就用 Padé 重求和
        \emph{推断真值}。优点：每步只需要一次稀疏 GEMV/GEMM——总成本
        \(\mathcal{O}(D^2)\) 而非 \(\mathcal{O}(D^3)\)，且可以
        \emph{只在 on-shell 行}做矩阵乘法；缺点：原始 Neumann 级数
        在三核能区内\textbf{常常发散}（\(\rConv < 1\)），
        必须靠 Padé 重求和才能恢复物理值。
\end{itemize}

\paragraph{为什么 Neumann 级数发散？}
原因有三层：
\begin{enumerate}
  \item \textbf{核耦合强}：三核动力学中 \(\Aker\) 的谱半径
        \(\rho(\Kker)\) 在共振或近似束缚态 \(q\) 处可达 \(\gtrsim 1\)。
        在数学上这就是几何级数 \(\sum r^n\) 在 \(|r| \ge 1\) 时发散。
  \item \textbf{近束缚态极点}：\((\Iker - \Kker)^{-1}\)
        视为复变量 \(z\) 的解析函数时，\(z = 1\) 与某个真实物理极点
        \(z_p\) 距离很近（\(|z_p| < 1\) 或紧贴单位圆），
        而 Neumann 级数的收敛半径 \(\rConv = |z_p|\)。
  \item \textbf{杂散幅度（spurious amplification）}：在 SWP 基下
        \(\Aker\) 是\emph{非对称}且\emph{非自伴}，特征值可能落到
        复平面单位圆外。
\end{enumerate}
正是这三条让 \emph{Padé 重求和成为求解 AGS 方程的物理必需品}，
而不是一种"加速优化"。Sean Miller 的博士论文与 PRC 106, 024001 (2022) 中
反复强调：\textbf{读者切勿把 Padé 当成一个"提速技巧"——它是把
形式发散的级数解析延拓到物理值的唯一手段}。

\paragraph{\texttt{solve\_dense} 的诊断价值。}
即便 \texttt{solve\_dense} 永远不能上生产规模，它在调试中无可替代：
当某个新势模型/新通道/新置换缓存怀疑出错时，先把 \(\Np=\Nq=10\) 跑一次
\texttt{solve\_dense=true}，记录 LAPACK 解出的"真"\(\Uker\)；
再以\emph{相同的输入}跑 Padé 路径，对比同一通道、同一 \(q\) 索引下的
\(\Uampl{\PWchanp}{\PWchan}\)。两者应当吻合到 \(\sim 10^{-7}\)（单精度极限），
吻合即认为\textbf{矩阵装配（\(\Aker, \Gker\)）正确}，
不吻合则锁定到 Padé 路径或装配路径其中之一。
commit \texttt{4d120d5}（"Removed max\_TFC as input, added solve\_dense.
solve\_dense solves the Faddeev equations using LAPACK, its slow and costly but
useful for debugging"）正是把这个 oracle 显式写进代码并固化下来。

\paragraph{读者应当带走什么。}
\begin{itemize}
  \item 看到一个 AGS 矩阵方程时，能在脑中立刻区分
        \(\Aker, \Kker = \Gker\Aker, \Pade{N}{N}\)
        三个层级；
  \item 给定 Neumann 系数序列 \(\{a_n\}_{n=0}^{2N}\) 能在纸上
        手写出 \(\Pade{N}{N}\) 的两条行列式公式；
  \item 在 solver log 里能一眼读出 "PA[N,M]"、"a\_n"、
        "Convergence reached for all on-shell elements" 三个标志的物理含义；
  \item 选 dense 还是 Padé 时知道阈值——\(D \lesssim 5000\) 用 dense，
        否则只能用 Padé；
  \item 看到 \(\Pade{N}{N}\) 跳跃式振荡时知道是\textbf{杂极点}
        而不是物理共振，会用 \texttt{convergence\_criteria\_2}
        的 \emph{minimum-PA-diff} 启发式判定真收敛阶。
\end{itemize}

% =========================================
\section{数学推导}
\label{sec:ch11-math}

% -----------------------------------------
\subsection{Neumann 级数的代数形式}
\label{subsec:ch11-neumann-derivation}

\paragraph{形式幂级数。}
\cref{eq:ch11-fixed-point} 的形式解是
\begin{equation}
\label{eq:ch11-neumann-formal}
\Uker \;=\; \sum_{n=0}^{\infty} \Kker^n\,\Aker
\;=\; \Aker + \Kker\Aker + \Kker^2\Aker + \cdots
\end{equation}
对应代码里的物理图像：\(n=0\) 项是\textbf{驱动项}（一次 \(\Vop\)
作用、不传播），\(n=1\) 项是"一次再散射"（\(\Vop\) 后传播一段、再 \(\Vop\)），
依此类推。把第 \(n\) 项的 on-shell 投影写成\textbf{标量系数}
\begin{equation}
\label{eq:ch11-neumann-scalar}
a_n^{(\PWchanp,\PWchan,q)} \;:=\;
\bigl[\Aker (\Gker \Aker)^n\bigr]_{(\PWchanp,q,p_d),\,(\PWchan,q,p_d)},
\end{equation}
其中 \(p_d\) 是 deuteron 束缚态在 SWP 谱中的固定索引（约定为 0）；
\(\PWchanp/\PWchan\) 是部分波通道。这个标量序列正是
\cref{lst:ch11-neumann-recursion} 代码中 \texttt{a\_coeff\_array} 与磁盘文件
\texttt{neumann\_terms\_*.txt} 存的内容。

\paragraph{收敛性判据。}
设 \(\Kker\) 的谱半径 \(\rho(\Kker) := \max_i |\lambda_i|\)，那么
\cref{eq:ch11-neumann-formal} 在 \(\rho(\Kker) < 1\) 时\emph{绝对收敛}；
反之级数发散。但即便发散，\((\Iker - \Kker)^{-1}\) 这个解析函数本身
\emph{依然存在}——只要 \(\det(\Iker - \Kker) \ne 0\)。这正是
Padé 重求和作为\textbf{解析延拓工具}的用武之地：把发散的级数
解析延拓回函数的"真值"。

% -----------------------------------------
\subsection{Padé 表 \texorpdfstring{$[N/M]$}{[N/M]} 的行列式公式}
\label{subsec:ch11-pade-table}

\paragraph{定义。}
给定一个标量级数 \(f(z) = \sum_{n=0}^{\infty} a_n z^n\)，其
Padé 表的 \([N/M]\) 元素是\emph{有理函数}
\begin{equation}
\label{eq:ch11-pade-def}
\Pade{N}{M}(z) \;=\; \frac{P_N(z)}{Q_M(z)},
\quad
P_N(z) = \sum_{k=0}^{N} p_k z^k,
\quad
Q_M(z) = \sum_{k=0}^{M} q_k z^k,
\end{equation}
满足 \emph{Padé 匹配条件}
\begin{equation}
\label{eq:ch11-pade-match}
f(z)\,Q_M(z) - P_N(z) \;=\; \mathcal{O}(z^{N+M+1}).
\end{equation}
即 \(\Pade{N}{M}\) 的 Taylor 展开前 \(N+M+1\) 阶与 \(f(z)\) 完全吻合。
求解 \(\{p_k, q_k\}\) 为线性代数问题：先解 \(M\) 个齐次方程
\(\sum_{j=0}^{M} q_j a_{N-M+1+i+j} = 0\)（\(i = 0, \ldots, M-1\)）确定
分母系数（约定 \(q_0 = 1\)），再代回求 \(p_k\)。

\paragraph{Heine--Frobenius 行列式形式。}
解上述线性系统的\textbf{闭式解}可写成两个 \((M+1)\times(M+1)\) 行列式之比：
\begin{equation}
\label{eq:ch11-pade-det}
\Pade{N}{M}(z)
\;=\;
\frac{
\det\!\begin{pmatrix}
\sum_{j=M}^{N} a_{j-M}\,z^j & \sum_{j=M-1}^{N} a_{j-M+1}\,z^j & \cdots & \sum_{j=0}^{N} a_{j}\,z^j \\
a_{N-M+1} & a_{N-M+2} & \cdots & a_{N+1} \\
a_{N-M+2} & a_{N-M+3} & \cdots & a_{N+2} \\
\vdots & \vdots & \ddots & \vdots \\
a_{N} & a_{N+1} & \cdots & a_{N+M}
\end{pmatrix}
}{
\det\!\begin{pmatrix}
z^M & z^{M-1} & \cdots & 1 \\
a_{N-M+1} & a_{N-M+2} & \cdots & a_{N+1} \\
\vdots & \vdots & \ddots & \vdots \\
a_{N} & a_{N+1} & \cdots & a_{N+M}
\end{pmatrix}
}.
\end{equation}
这正是 \cref{lst:ch11-pade-approximant} 里函数 \texttt{pade\_approximant} 的字面落地：
两个 \((M+1)\times(M+1)\) 复数矩阵 \texttt{P\_array}/\texttt{Q\_array} 各自调用一次
\texttt{determinant()} 后相除。

\paragraph{对角 Padé \texorpdfstring{$\Pade{N}{N}$}{[N/N]} 的特殊地位。}
本仓库\emph{固定使用对角元素 \(\Pade{N}{N}\)}（\texttt{NM\_max=14}）。
原因：
\begin{enumerate}
  \item 对角 Padé 是\emph{所有} \([N/M]\) 中收敛域\textbf{最大}的子序列
        （Stahl 极小化定理）；
  \item 对角 Padé 对系数 \(a_n\) 的舍入误差最稳定（条件数 \(\sim N\) 而非 \(\sim N^2\)）；
  \item \(\Pade{N}{N}\) 用 \(2N+1\) 个 Neumann 系数；这正好让
        \cref{eq:ch11-fixed-point} 的迭代步数 \(n=0, 1, 2, \ldots, 2N\)
        与 Padé 阶数 \(N=0, 1, \ldots, N_{\max}\) 同步\textbf{一对一}增长。
\end{enumerate}
具体地，从 \(\Pade{N-1}{N-1}\) 升级到 \(\Pade{N}{N}\) 只需要再生成
\(a_{2N-1}\) 与 \(a_{2N}\) 两个新系数（即代码循环 \texttt{for n=2*NM-1; n<2*NM+1}），
再调一次行列式求 \(\Pade{N}{N}(z=1)\)。

\paragraph{计算点 \texorpdfstring{$z=1$}{z=1}：把级数当几何级数。}
本仓库统一在 \(z=1\) 处求值：\(\Pade{N}{N}(1)\) 即近似 \(\sum_n \Kker^n \Aker\)
形式作用到入射态后的散射振幅。这等价于把 Neumann 级数当成在
\(z=1\) 处展开的形式幂级数：原级数 \(\sum a_n\) 等价于
\(f(1) = \sum a_n \cdot 1^n\)，即使在 \(\rConv < 1\) 时也可由 \(\Pade{N}{N}(1)\)
进行\emph{解析延拓}。这跟物理图像吻合：\(z=1\) 表示
"完整的 \(\Gker\) 与 \(\Aker\) 都施加一次"。

\paragraph{Wynn 的 \(\eps\) 算法（背景）。}
教科书上 Padé 表的另一种生成法是 Wynn \(\eps\) 算法，
通过递推 \(\eps_{k+1}^{(n)} = \eps_{k-1}^{(n+1)} + 1/(\eps_k^{(n+1)} - \eps_k^{(n)})\)
搭出整张表。本仓库\textbf{未}使用 Wynn——它在 \(\Aker\) 复数化后的
数值稳定性比直接行列式公式更差（除以小差分项），且无法直接给出
对角元的代数表达。Sean 路径选择闭式行列式：好处是 14 阶之内不会出
"\(0/0\)" 数值病态，坏处是每次升阶都要重算 \((M+1)^2\) 个矩阵元的复行列式。

% -----------------------------------------
\subsection{收敛半径与极点分布}
\label{subsec:ch11-radius-pole}

\paragraph{Borel--Padé 视角。}
若 \(f(z)\) 形式幂级数发散，但其\textbf{Borel 变换}
\(B(z) = \sum_n a_n z^n / n!\) 收敛，则 \(f(z) = \int_0^{\infty} e^{-t} B(zt)\,dt\)
给出\emph{解析延拓}（Borel 求和）。Padé 表 \(\Pade{N}{N}\) 在很多情形下
\textbf{隐式实现 Borel 求和}：当 \(\Pade{N}{N}(z)\) 的极点逼近真实 \(f(z)\)
的 Riemann 面奇点时，序列 \(\{\Pade{N}{N}(z=1)\}_{N=0}^{\infty}\) 仍然收敛。
这是 Padé 在 AGS 求解中"明知道发散仍能用"的根本原因。

\paragraph{Stahl 收敛域。}
设 \(f(z)\) 在复平面去掉某条切割集后解析（典型：
两体阈值切线 + 三体阈值切线 + 离散极点）。Stahl 定理保证
\(\Pade{N}{N}(z) \to f(z)\) 在该切割集的\emph{补集}上一致收敛。
对于三体 \(d+p\) 散射在 \(\Tlab \sim 100\sim 250\) MeV，最近的奇点是
\((d, p) \to (n, p, p)\) 阈值与 deuteron 极点；Stahl 域包含
\(z=1\) 这一物理求值点，故 \(\Pade{N}{N}(1)\) 必然收敛——
事实上 \cref{sec:ch11-numerical} 中给出的数值显示 \(N=4\sim 6\) 已经收敛到 6 位有效数字。

% -----------------------------------------
\subsection{Kowalski--Noyes 形式与对照}
\label{subsec:ch11-kn-form}

Kowalski--Noyes (KN) 把 \cref{eq:ch11-fixed-point} 改写为
\(\Uker = (\Iker - \Kker)^{-1}\Aker\) 后引入\emph{半 on-shell 投影矩阵}
\(\Pi\)：\(\Uker_{\text{HOS}} = \Pi\Uker\Pi\)。WPCD 路径下，
"\(\Pi\)"是把第 \(p\) 维约束到 \(p=p_d\)（束缚态 SWP 索引）的投影；
KN 求解器存的就是这个 \(\NqWP \times \NqWP\) 半 on-shell 子块。

WPCD/Padé 路径与 KN 的关系：两者解的是\emph{同一个}方程，但
WPCD 把"解 \(\Uker\) 的整块"转化为"求 on-shell 元素的标量级数 \(\{a_n\}\) 然后 Padé 重求和"，
省去了 \(\Pi\) 投影矩阵的显式存储；代价是必须每个 \(q_{\text{shell}}\)
独立做一次 \(\sim 2\,N_{\max}+1\) 步的 Neumann 推进。
KN 路径计算量 \(\mathcal{O}(D \cdot N_{q})\)，
WPCD/Padé 路径 \(\mathcal{O}(D \cdot N_{q} \cdot N_{\max})\)；
但后者每步矩阵乘法是 BLAS-3 GEMM，常数因子小得多。

% -----------------------------------------
\subsection{Dense 路径的 \texorpdfstring{$(I-K)$}{(I-K)} 直接求逆}
\label{subsec:ch11-dense-derivation}

\paragraph{矩阵方程改写。}
\cref{eq:ch11-fixed-point} 等价于 \emph{矩阵} \(\Uker\) 满足
\(\Lker\Uker = \Rker\)，其中
\begin{equation}
\label{eq:ch11-LU-RHS}
\Lker = \Iker - \Gker\Aker, \qquad \Rker = \Aker.
\end{equation}
注意 \cref{lst:ch11-dense-solver} 里 \(\Aker\) 由 \texttt{calculate\_CPVC\_col} 一列一列生成
（不再先组合成稀疏 \(\Pker\Vker\Cker\)，而是 \emph{每列即时}计算），
然后 \(\Lker\) 块按 \(L_{\rho \sigma} = -A_{\rho \sigma} G_\sigma\) 装配，
对角加 \(+1\) 完成 \(\Iker - \Gker\Aker\)。最后调用
\texttt{solve\_MM(L, R, dense\_dim)}（封装 LAPACK \texttt{ZGESV}）
解出 \(\Rker \leftarrow \Lker^{-1}\Rker\)。

\paragraph{LAPACK 内核：\texttt{ZGESV}。}
\texttt{ZGESV}（复双精度稠密 LU 求解）做：
(1) \(\Lker = P\,L\,U\) 部分主元 LU 分解，\(\mathcal{O}(D^3/3)\)；
(2) \(\Lker\,X = \Rker\) 通过两步三角求解 \(L Y = P\Rker\)，
\(U X = Y\)，\(\mathcal{O}(D^2 \cdot N_{\text{rhs}})\)。
对于 dense 路径 \(N_{\text{rhs}} = D\)（解 \(D\) 列），故总
\(\mathcal{O}(D^3)\)。\(\Aker\) 只在 \(\Pker\) 可缓存时减半工作量。

\paragraph{为什么 dense 路径写出来。}
\begin{itemize}
  \item \textbf{benchmark oracle}：toy-size \(\Np=\Nq=10\) 跑 dense
        给出"无迭代误差"的参考振幅，对照 Padé 路径定位 bug；
  \item \textbf{诊断中间矩阵}：dense 路径每个 \(q\) 都把
        \(\Aker, \Lker, \Uker\) 三个 \(D\times D\) 矩阵写到磁盘
        （\texttt{A\_array\_E\_idx\_*.txt} 等），方便 Python 后处理；
  \item \textbf{杂极点排查}：当怀疑 Padé 路径出杂极点（\cref{sec:ch11-pitfalls}），
        把 dense 路径的 \(\Lker\) 拿出来直接看
        \(\det(\Lker)\) 是否过小或异号。
\end{itemize}

\subsection{矩阵装配的代价模型}
\label{subsec:ch11-cost}

把 \(\Aker\) 视为稠密：内存 \(\mathcal{O}(D^2)\)，
其中 \(D = \Nalpha\,\NqWP\,\NpWP\)。例如典型生产规模
\(\Nalpha = 64, \NqWP = 30, \NpWP = 30\) 时
\(D \approx 5.8 \times 10^4\)，\(D^2\) 个 \texttt{double}
就是 \(\sim\!27\) GiB——\textbf{这是 dense 路径在生产规模下不可行的根本原因}。
但 \(\Pker\) 是稀疏的（密度 \(\sim 0.1\%\)），\(\Aker = \Cker^\top \Pker \Vker \Cker\)
的"稀疏 \(\Pker\) 主导"特性使得每\emph{一列}的稀疏度只有 \(\sim D \times 10^{-3}\)。
Padé 路径正是利用这一点：从未材料化整块 \(\Aker\)，而是
\emph{按需逐列重生成}（\(D\) 个列 \(\to\) \(D \cdot D \cdot 10^{-3} = D^2 \cdot 10^{-3}\)
存储）。再加上"压缩到非收敛行"的策略，实测内存只需 \(\sim 4\) GiB。

% =========================================
\section{离散化}
\label{sec:ch11-discretization}

\paragraph{在 SWP 基下写出所有矩阵。}
本节明确写出本章后续所用矩阵的 SWP-基定义，方便交叉引用。

\subsection{驱动矩阵 \texorpdfstring{$\Aker$}{A} 的指标结构}
\label{subsec:ch11-A-index}
在 SWP 基（\cref{ch:06_wp_swp_states}）下
\begin{equation}
\label{eq:ch11-A-def}
A_{(\PWchanp i' j'),(\PWchan i j)}
\;=\;
\sum_{\PWchanp', \PWchan'}
\bigl[\Cker^\top\bigr]_{\PWchanp \PWchanp'}\,
\bigl[\Pker\bigr]_{(\PWchanp' i' j'),(\PWchan' i j)}\,
\bigl[\Vker\bigr]_{\PWchan' \PWchan}\,
\bigl[\Cker\bigr]_{\PWchan \PWchan},
\end{equation}
其中 \(i\) 是 \(p\)-bin 索引、\(j\) 是 \(q\)-bin 索引。
代码里把这一三重指标 \((\PWchan, j, i)\) 通过 \texttt{dense\_packet\_index}
（\cref{lst:ch11-row-compaction-helper} 上方）打包成单一 \texttt{size\_t}：
\begin{equation}
\label{eq:ch11-flat-idx}
\sigma \;=\; \PWchan \cdot \NqWP \cdot \NpWP + j\cdot\NpWP + i,
\qquad
\sigma \in [0, D),\quad D = \Nalpha\,\NqWP\,\NpWP.
\end{equation}
SWP 基下 \(\Vker\) 块对角（每个 \(\PWchan\) 内只在自己的 LS 子空间中），
\(\Pker\) 高度稀疏（密度 \(\sim 0.1\%\)），\(\Cker\) 几乎对角（\(\Np\times\Np\) 块），
\(\Gker\) 在 SWP 基下\emph{严格对角}——这是
\cref{ch:10_resolvent} 把 \(\Gop_0\) 拆成 \(R + Q\)
形式的几何意义。

\subsection{Neumann 级数的矩阵幂展开}
\label{subsec:ch11-neumann-mat-expansion}
代入 \(\Kker = \Gker\Aker\) 后，\cref{eq:ch11-neumann-formal}
按矩阵幂展开为
\begin{equation}
\label{eq:ch11-neumann-mat}
U_{\rho\sigma}
\;=\;
A_{\rho\sigma}
+ G_\rho A_{\rho\sigma_1} G_{\sigma_1} A_{\sigma_1 \sigma}
+ G_\rho A_{\rho\sigma_1} G_{\sigma_1} A_{\sigma_1\sigma_2} G_{\sigma_2} A_{\sigma_2\sigma}
+ \cdots
\end{equation}
（求和 \(\sigma_1, \sigma_2, \ldots\) 隐含；\(G_\rho\) 是对角元）。
注意每相邻的 \(A\) 之间夹一个对角的 \(G\)；\cref{lst:ch11-neumann-recursion}
里的 \emph{Multiply An by G} 块（"复数 \(\Gker\) 元素的 4 个实数乘法"）正是这一项。

\subsection{on-shell 子集的提取}
\label{subsec:ch11-on-shell-extract}
最终我们只需要 \(U_{\rho\sigma}\) 的极少数 on-shell 元素：
\begin{equation}
\label{eq:ch11-on-shell-subset}
\Uampl{\PWchanp}{\PWchan}(q)
\;=\;
U_{(\PWchanp,\, j_q,\, p_d),\,(\PWchan,\, j_q,\, p_d)}.
\end{equation}
其中 \(j_q\) 是给定 \(\Tlab\) 对应的 on-shell \(q\)-bin（\texttt{q\_com\_idx\_array}），
\(p_d\) 是 deuteron 在 SWP 谱中的 bin（约定固定 0；\cref{ch:06_wp_swp_states} §6.4）。
本仓库共 \(N_{\text{NDOS}} = N_{d-\text{row}} \cdot N_{d-\text{col}} \cdot N_{q-\text{com}}\)
个 on-shell 元素，每个独立做 Neumann+Padé；这就是 \texttt{a\_coeff\_array}
按 \texttt{idx\_NDOS} 分块组织的原因。

% =========================================
\section{算法骨架}
\label{sec:ch11-algorithms}

% --- Algo 1 ---
\begin{algorithm}[H]
\caption{NEUMANN-ITER（Neumann 级数生成器，对应 \texttt{pade\_method\_solve} 主循环）}
\label{alg:ch11-neumann-iter}
\KwIn{\(\Aker\)（按列构造）, \(\Gker\)（对角，每个 \(q\) 独立）, on-shell 行索引集 \(\mathcal{R}\), \(N_{\max}\)}
\KwOut{Neumann 系数 \(\{a_n^{(\rho,\sigma)}\}_{n=0}^{2N_{\max}}\)}
\BlankLine
\(\Aker_0 \leftarrow \Aker\) \tcp*{即 \texttt{re/im\_A\_An\_row\_array\_prev}}
\For{\(\rho \in \mathcal{R}\), \(\sigma \in\mathcal{S}_{\text{NDOS}}\)}{
  \(a_0^{(\rho,\sigma)} \leftarrow [\Aker_0]_{\rho\sigma}\)\;
}
\For{\(n = 1\) \KwTo \(2N_{\max}\)}{
  \tcc{Step 1: Multiply An by G (diagonal)}
  \For{\(\rho\in\mathcal{R}\), \(i \in [0,D)\)}{
    \(a^{\text{prev}}_{\rho i} \leftarrow [\Aker_{n-1}]_{\rho i}\,G_i\)\;
  }
  \tcc{Step 2: row compaction (skip converged rows)}
  \(\mathcal{R}_{\text{active}} \leftarrow \{\rho \in \mathcal{R} : \text{has unconverged target}\}\)\;
  \tcc{Step 3: GEMM with CPVC, regenerated in chunks}
  \For{每个列 chunk \(\mathcal{C}\)}{
    生成 \texttt{CPVC\_cols\_array} \(\leftarrow \Aker[:,\mathcal{C}]\)\;
    \(\Aker_n[\mathcal{R}_{\text{active}}, \mathcal{C}] \leftarrow \Aker_{n-1}[\mathcal{R}_{\text{active}}, :] \cdot \Aker[:,\mathcal{C}]^\top\)
    \tcp*{cblas\_dgemm 实复各一次}
  }
  \tcc{Step 4: 提取 on-shell 标量}
  \For{每个 NDOS \((\PWchanp,\PWchan,q)\)}{
    \(a_n^{(\PWchanp,\PWchan,q)} \leftarrow [\Aker_n]_{(\PWchanp,j_q,p_d),(\PWchan,j_q,p_d)}\)\;
  }
}
\end{algorithm}

% --- Algo 2 ---
\begin{algorithm}[H]
\caption{PADE-RESUM（对角 Padé \(\Pade{N}{N}\) 计算与收敛判定）}
\label{alg:ch11-pade-resum}
\KwIn{标量级数 \(\{a_n\}_{n=0}^{2N_{\max}}\), \(N_{\max}\)}
\KwOut{最佳 \(\Pade{N^*}{N^*}(z=1)\)}
\BlankLine
\For{\(N = 0\) \KwTo \(N_{\max}\)}{
  生成 \((M+1)\times(M+1)\) 矩阵 \(P_{\text{arr}}, Q_{\text{arr}}\) 按 \cref{eq:ch11-pade-det}\;
  \(\Pade{N}{N} \leftarrow \det(P_{\text{arr}}) / \det(Q_{\text{arr}})\)\;
  \(\Delta_N \leftarrow |\Pade{N}{N} - \Pade{N-1}{N-1}|\)\;
  \tcc{Convergence criteria (任一满足即认为收敛)}
  \If{\(N = N_{\max}\)}{break\tcp*{Criterion 0}}
  \If{\(N - \arg\min_{N'} \Delta_{N'} > 4\)}{break\tcp*{Criterion 1}}
  \If{\(\min_{N'}\Delta_{N'} < 10^{-6}\,|\Pade{N^*}{N^*}|\)}{break\tcp*{Criterion 2}}
  \If{\(\min_{N'}\Delta_{N'} < 10^{-7}\)}{break\tcp*{Criterion 3}}
}
\Return \(\Pade{N^*}{N^*}\) where \(N^* = \arg\min_{N'} \Delta_{N'}\)\;
\end{algorithm}

% --- Algo 3 ---
\begin{algorithm}[H]
\caption{DENSE-SOLVE（基准用稠密 LAPACK 求解器，对应 \texttt{faddeev\_dense\_solver}）}
\label{alg:ch11-dense-solve}
\KwIn{\(\Aker, \Gker\), on-shell 元素列表}
\KwOut{所有 \(\Uker\) 元素（含 off-shell；后处理只用 on-shell）}
\BlankLine
\For{每个 \(q\) shell}{
  分配 \(\Lker, \Rker \in \mathbb{C}^{D\times D}\)\;
  \For{每列 \(\sigma = (\PWchan, j, i)\)}{
    \(\Aker[:,\sigma] \leftarrow\) \texttt{calculate\_CPVC\_col}(...)\;
    \(\Lker[:,\sigma] \leftarrow -\Aker[:,\sigma]\,G_\sigma\), \(\Rker[:,\sigma] \leftarrow \Aker[:,\sigma]\)\;
  }
  \(\Lker[\sigma,\sigma] \mathrel{+}= 1\)\tcp*{即 \(\Iker - \Gker\Aker\)}
  \texttt{solve\_MM}(\(\Lker, \Rker, D\))\tcp*{ZGESV: LU + TRS}
  从 \(\Rker\) 中按 NDOS 索引拷贝到 \texttt{U\_array}\;
}
\end{algorithm}

% --- Algo 4 ---
\begin{algorithm}[H]
\caption{CONVERGENCE-DIAG（收敛诊断与压缩行表）}
\label{alg:ch11-conv-diag}
\KwIn{\texttt{pade\_approximants\_conv\_array}, \texttt{pade\_approximants\_BU\_conv\_array}}
\KwOut{活动行集合 \(\mathcal{R}_{\text{active}}\), 全收敛标志}
\BlankLine
\(\mathcal{R}_{\text{active}} \leftarrow \emptyset\)\;
\For{每个 (deuteron-row \(d\), q-shell \(q\))}{
  \tcc{check whether all NDOS targeting (d, q) have converged}
  \(\text{all\_done} \leftarrow \bigwedge_{d'} \text{conv}[d, d', q]\)\;
  \If{include\_breakup}{
    \(\text{all\_done} \mathrel{\&}= \bigwedge_{\text{BU chn}} \text{BU\_conv}[d, \cdot]\)\;
  }
  \If{not all\_done}{
    \(\mathcal{R}_{\text{active}} \leftarrow \mathcal{R}_{\text{active}} \cup \{(d, q)\}\)\;
  }
}
\Return \(\mathcal{R}_{\text{active}}\), \(|\mathcal{R}_{\text{active}}| = 0\)\;
\end{algorithm}

\paragraph{算法之间的耦合。}
NEUMANN-ITER 与 PADE-RESUM 是\textbf{交错执行}的：
每生成 2 个新 Neumann 项就升一阶 Padé（\(N\to N+1\)），然后调用
CONVERGENCE-DIAG 判定哪些行已经收敛、可以压缩。这一交错设计来自
commit \texttt{cfaa081}（"Rewrote ordering of loops in faddeev solver to
calculate pade approximants for each new couple of Neumann terms, and check
convergence. This adds a noticable optimisation rather than calculating all
neumann terms first"）；它把"早收敛行"
的后续 GEMM 工作量直接砍掉。DENSE-SOLVE 则与三者无交错——
它是独立分支，不参与 Neumann/Padé 交错。

% =========================================
\section{代码落地}
\label{sec:ch11-code}

% -----------------------------------------
\subsection{主入口签名：\texttt{solve\_faddeev\_equations}}
\label{subsec:ch11-entry}

\cref{ch:04_faddeev_ags} 已经引用过同函数的\textbf{头文件签名}
（\texttt{solve\_faddeev.h}）。这里我们抽 \texttt{solve\_faddeev.cpp}
里的\textbf{实现起手段}（locals 解包），它揭示了下游 Padé/dense 分支
\emph{共用}的预处理：

\lstinputlisting[style=cpp,
  caption={current 入口实现起手段（\texttt{solve\_faddeev.cpp}, 1825--1850 行）},
  label={lst:ch11-entry-current}]
  {code_excerpts/snippets/ch11_solve_faddeev_entry.cpp}

origin 仓库的同函数签名\textbf{无 \texttt{tnf} 参数}（3NF 还未接入）：

\lstinputlisting[style=cpp,
  caption={origin 入口签名（\texttt{CPP/solve\_faddeev.cpp}, 1446--1466 行；与 current 对照）},
  label={lst:ch11-entry-origin}]
  {code_excerpts/snippets/ch11_solve_faddeev_entry_origin.cpp}

\paragraph{对照要点。}
\begin{itemize}
  \item current 在末尾增加 \texttt{const three\_nucleon\_force\_model* tnf = nullptr}，
        并构造 \texttt{tnf\_kernel\_context}（默认值 null 表示 3NF 关闭，
        热路径仅一次 \texttt{enabled()} 测试即跳过——与 \cref{ch:15_3nf_physics} 衔接）；
  \item current 的参数 \texttt{fwp\_statespace fwp\_states} 是新加的，
        提供 \texttt{p\_WP\_array, q\_WP\_array} 给 3NF 核访问 WP 边界；
  \item 局部变量解包逻辑两版完全一致——这一段是\textbf{接口稳定性}的
        证据：3NF 引入只在签名末尾追加，不破坏 origin 的其他调用位点。
\end{itemize}

% -----------------------------------------
\subsection{Padé 重求和的核心：\texttt{pade\_approximant}}
\label{subsec:ch11-pade-approximant}

\lstinputlisting[style=cpp,
  caption={\texttt{pade\_approximant}：行列式比值（\texttt{solve\_faddeev.cpp}, 649--674 行）},
  label={lst:ch11-pade-approximant}]
  {code_excerpts/snippets/ch11_pade_approximant.cpp}

\paragraph{逐行对照 \cref{eq:ch11-pade-det}。}
\begin{itemize}
  \item \texttt{P\_array} 与 \texttt{Q\_array} 都是 \((M+1)^2\) 复数；
        前 \(M\) 行是\textbf{完全相同}的"系数 Hankel 块"
        \(a_{N-M+1+i+j}\)（\(i = 0,\ldots,M-1\)；\(j = 0,\ldots,M\)），
        对应行列式公式中下面 \(M\) 行；
  \item 第 \(M\) 行（最后一行）不同：\texttt{Q\_array} 填 \(z^{M-k}\)，
        \texttt{P\_array} 填\textbf{部分和} \(\sum_{j=M-k}^{N} a_{j-(M-k)} z^j\)；
        这正对应 \cref{eq:ch11-pade-det} 的最上行/底行；
  \item 最后两个 \texttt{determinant()} 调用与一次复数除法即得 \(\Pade{N}{M}(z)\)。
\end{itemize}
注意 \texttt{determinant} 来自 \texttt{utils/matrix\_routines}，对 \(M\le 14\)
是直接 LU 分解；对更大 \(M\) 会出现行列式数值溢出/下溢——
这是 \texttt{NM\_max=14} 这一 hard cap 的原因。

\paragraph{origin 与 current 几乎逐字符相同。}
该函数从 \texttt{161c557}（"Pade method works, benchmarked with sparse-matrix
python code"）一直没改过；origin 与 current 仅缩进微差。

\lstinputlisting[style=cpp,
  caption={origin 同函数（\texttt{CPP/solve\_faddeev.cpp}, 271--296 行）作为对照},
  label={lst:ch11-pade-approximant-origin}]
  {code_excerpts/snippets/ch11_pade_approximant_origin.cpp}

% -----------------------------------------
\subsection{Neumann 递推与 GEMM 主循环}
\label{subsec:ch11-neumann-recursion}

\lstinputlisting[style=cpp,
  caption={Neumann 递推主循环外层（\texttt{solve\_faddeev.cpp}, 1304--1395 行）},
  label={lst:ch11-neumann-recursion}]
  {code_excerpts/snippets/ch11_neumann_recursion.cpp}

\paragraph{走读这段代码。}
\begin{enumerate}
  \item \texttt{for (size\_t NM=0; NM<NM\_max+1; NM++)}：外层是 Padé 阶；
        每次升一阶；若 \texttt{num\_converged\_elements==num\_EL\_A\_vals}
        则全收敛 \texttt{break}（对应 \cref{alg:ch11-conv-diag}）；
  \item 内层 \texttt{for (int n=2*NM-1; n<2*NM+1; n++)}：每升一阶 Padé
        生成 2 个新 Neumann 项（\(n=2N-1\) 与 \(n=2N\)）；这是
        commit \texttt{cfaa081} 引入的"\texttt{couple of Neumann terms}"风格；
  \item \texttt{Multiply An by G}：对角矩阵-密集向量乘，复数实部虚部交叉，
        4 个实数乘法对应 1 个复数乘法（\cref{eq:ch11-neumann-mat}
        中相邻 \(A\) 间的 \(G\)）；
  \item \texttt{row\_has\_only\_converged\_targets} 行压缩：
        把已收敛行从 \texttt{prev} 拷到 \texttt{comp}，活动行数
        \texttt{num\_non\_conv\_rows}；
  \item 之后的 \texttt{cblas\_dgemm} 才是真正的工作量——
        见接下来的 GEMM 片段。
\end{enumerate}

\lstinputlisting[style=cpp,
  caption={Neumann 递推核心 GEMM（\texttt{solve\_faddeev.cpp}, 1444--1505 行）},
  label={lst:ch11-neumann-gemm}]
  {code_excerpts/snippets/ch11_neumann_gemm.cpp}

\paragraph{为什么是两次 \texttt{dgemm} 而不是一次 \texttt{zgemm}？}
本仓库\textbf{把复矩阵拆成实部/虚部两个独立 \texttt{dgemm}}：
\begin{equation}
\label{eq:ch11-complex-as-real}
(\re A_n + i\,\im A_n)\,\Aker
\;=\;
\re A_n\,\Aker + i\,\im A_n\,\Aker.
\end{equation}
这成立的前提是 \(\Aker\)（\texttt{CPVC\_cols\_array}）\textbf{严格实}——
正是本仓库的真实情形（\(\Cker, \Pker, \Vker\) 都是实矩阵；
\(\Gker\) 是复，但只在 \texttt{Multiply An by G} 一步施加）。
两次 \texttt{dgemm} 比一次 \texttt{zgemm} 快约 1.7 倍：因为 \texttt{zgemm}
把每个矩阵元当 16 字节复数处理，BLAS 内部要做 4 次实数乘和 2 次实数加；
而我们已经把 \(\Aker\) 实数化，只剩 1 次实数乘。
若将来引入复数 NN 势（如 optical model），需改回 \texttt{zgemm}。

% -----------------------------------------
\subsection{行压缩 helper：\texttt{row\_has\_only\_converged\_targets}}
\label{subsec:ch11-row-compaction}

\lstinputlisting[style=cpp,
  caption={收敛行检测 helper（\texttt{solve\_faddeev.cpp}, 86--122 行；anonymous namespace）},
  label={lst:ch11-row-compaction-helper}]
  {code_excerpts/snippets/ch11_row_compaction_helper.cpp}

\paragraph{物理含义。}
这个 helper 是 commit \texttt{30338ac}（"Improved Neumann-series convergence
data printout"）引入的可读性重构产物——origin 把整段判定写在嵌套 \texttt{if}
里、混在主循环中；current 提取成自由函数，逻辑\emph{完全等价}但意图清晰：
"当且仅当某个 (deuteron-row, q-shell) 对应的所有 (d-col)
弹性 NDOS \emph{以及}（若启用 BU）所有 BU 通道都已收敛，
该行才能从 \texttt{re/im\_A\_An\_row\_array\_comp} 中剔除。"
注意它是\emph{保守的}：只要有一个目标尚未收敛就保留整行——
这避免了"行已删除但还需要"的潜在 bug。

% -----------------------------------------
\subsection{Padé 收敛诊断与判据}
\label{subsec:ch11-pade-convergence}

\lstinputlisting[style=cpp,
  caption={Padé 收敛判据（\texttt{solve\_faddeev.cpp}, 1609--1679 行）},
  label={lst:ch11-pade-convergence}]
  {code_excerpts/snippets/ch11_pade_convergence.cpp}

\paragraph{4 个收敛判据。}
代码里 4 个 \texttt{convergence\_criteria\_k} 对应不同物理情境：
\begin{description}
  \item[\texttt{criterion 0}: \(N = N_{\max}\)] 顶到 \texttt{NM\_max=14}，
        无法再升阶——只能取当前最佳值（哪怕没"真"收敛）。这条触发常意味着
        Padé 在该通道根本没收敛；应当作为\textbf{warning}出现在 log 里
        （commit \texttt{30338ac} 的"convergence data printout"改进就指它）。
  \item[\texttt{criterion 1}: \(N - N^* > 4\)] 过去 4 阶都不如某个早期阶——
        说明已经"路过最佳点"，再升只会更差（杂极点开始主导）。
        这是\emph{stagnation 检测}。
  \item[\texttt{criterion 2}: \(\Delta < 10^{-6}|\Pade{N^*}{N^*}|\)] 相对稳定 6 位有效数字。
        实际生产中绝大多数收敛走这一条。
  \item[\texttt{criterion 3}: \(\Delta < 10^{-7}\)] 单精度极限。
        若 \(|\Pade{N^*}{N^*}| < 10^{-1}\)，criterion 2 不会触发，靠这个
        绝对阈值兜底。
\end{description}
\textbf{注意细节}：判据用的是 \(\Delta = |\Pade{N}{N} - \Pade{N-1}{N-1}|\)，
\textbf{不是}相对差；这意味着大幅振幅的杂极点（\(\sim 10^{12}\)）仍可能
误触发 criterion 2——这正是 \cref{sec:ch11-pitfalls} 中提到的杂极点风险。

\paragraph{\texttt{idx\_best\_PA} 选取。}
这段代码维护一个\textbf{滑动最小差分追踪}：从 \(N=0\) 到当前 \(N\) 中，
找到 \(\Delta_{N'}\) 最小那个 \(N^*\)，并把对应的
\texttt{pade\_approximants\_array[idx\_NDOS*(NM\_max+1) + N\^*]}
作为最终输出。这一启发式来自 commit \texttt{88378e2}
（"Some speedup to CPVC generation. 2 additional convergence criteria for
Pade approximant ..."）——具体说就是 criterion 2/3 的引入。

% -----------------------------------------
\subsection{Dense LAPACK 求解器：\texttt{faddeev\_dense\_solver}}
\label{subsec:ch11-dense-solver}

\lstinputlisting[style=cpp,
  caption={\texttt{faddeev\_dense\_solver} 主体（\texttt{solve\_faddeev.cpp}, 676--776 行）},
  label={lst:ch11-dense-solver}]
  {code_excerpts/snippets/ch11_dense_solver.cpp}

\paragraph{对照 \cref{alg:ch11-dense-solve}。}
\begin{itemize}
  \item \texttt{L\_array} 与 \texttt{R\_array} 是
        \(\mathcal{O}(D^2)\) 复双精度大块——
        典型 toy 规模 \(D \sim 1\,200\) 时单块 \(\sim 23\) MiB；
  \item 三重循环 \texttt{idx\_q\_c, idx\_alpha\_c, idx\_p\_c}
        把每一列 \(\sigma\) 对应的 CPVC-列计算进 \texttt{CPVC\_col\_array}，
        然后乘以 \(G_\sigma = \Gker[j*D + \sigma]\) 得到 \(\Lker\) 列；
  \item \texttt{L\_array[col\_idx*dense\_dim + col\_idx] += 1}：对角加 \(\Iker\)；
  \item 三个 \texttt{store\_array}：\textbf{诊断转储}。dense 路径专门把
        \(\Aker (= \Rker)\)、\(\Lker (= K)\)、求逆后的 \(\Uker\) 写文件，
        给 Python 后处理对照。这是 commit \texttt{4d120d5} 的关键贡献：
        "solve\_dense ... useful for debugging."。
  \item \texttt{solve\_MM(L\_array, R\_array, dense\_dim)}：本仓库
        在 \texttt{utils/matrix\_routines.h} 中封装 LAPACK \texttt{ZGESV}
        ——\(\Lker \to LU\)，再两步三角求解 \(\Lker X = \Rker\)，结果原地写回 \(\Rker\)。
\end{itemize}

\paragraph{为什么 origin 几乎逐字符相同。}
对照原 \texttt{CPP/solve\_faddeev.cpp} 305--438 行——除了
\texttt{make\_elastic\_on\_shell\_index} 重构之外，
\(\Lker, \Rker\) 装配与 \texttt{solve\_MM} 调用一字不差。
这反映 dense 路径在 commit \texttt{4d120d5} 写定后\textbf{基本冻结}——
作为 oracle 它不需要演进，演进的全在 Padé 路径。

% =========================================
\section{数字闭环}
\label{sec:ch11-numerical}

\begin{numericalbox}
\textbf{场景：\(\Tlab=190\) MeV 在 \(J^\pi=1/2^+\) 通道下的 Padé 收敛序列。}

\textbf{复现命令}（约 5 分钟单 J 单 \(T_{\text{lab}}\)）：
\begin{verbatim}
cd Tic-tac
python examples/deuteron_proton_Ay.py \
    --work-dir output/deuteron_proton_Ay_tlab190 \
    --target-tlab-mev 190 --reuse-p123
\end{verbatim}
其中 \(\Np=\Nq=20\)，\texttt{two\_J\_3N\_max=1}，\texttt{J\_2N\_max=2}，
N2LOopt 势。on-shell q-bin 索引 15、16（对应到达 \(\Tlab=190\)
\(\pm \mathcal{O}(10\,\text{MeV})\)）。

\textbf{Neumann 系数序列}（取 \(\PWchanp = \PWchan = 4\)、\(q=15\) 的弹性元素，
real + imag 两列；从 \texttt{output/deuteron\_proton\_Ay\_tlab190/solver\_out/neumann\_terms\_Np\_20\_Nq\_20\_JP\_1\_1\_Jmax\_2.txt}）：

\medskip
\begin{tabular}{rrr}
\toprule
\(n\) & \(\re a_n\) & \(\im a_n\) \\
\midrule
0  & \(-6.27\times 10^{-1}\) & \(0.00\) \\
1  & \(-1.18\times 10^{0}\)  & \(-5.72\times 10^{0}\) \\
2  & \(+2.84\times 10^{0}\)  & \(+6.06\times 10^{-1}\) \\
3  & \(+1.52\times 10^{-1}\) & \(+3.24\times 10^{0}\) \\
4  & \(-1.92\times 10^{0}\)  & \(+4.86\times 10^{-1}\) \\
5  & \(-9.28\times 10^{-1}\) & \(-2.43\times 10^{0}\) \\
6  & \(+1.93\times 10^{0}\)  & \(-9.33\times 10^{-1}\) \\
7  & \(+9.76\times 10^{-1}\) & \(+1.84\times 10^{0}\) \\
8  & \(-1.50\times 10^{0}\)  & \(+1.02\times 10^{0}\) \\
9  & \(-1.07\times 10^{0}\)  & \(-1.33\times 10^{0}\) \\
10 & \(+1.11\times 10^{0}\)  & \(-1.08\times 10^{0}\) \\
11 & \(+1.07\times 10^{0}\)  & \(+9.39\times 10^{-1}\) \\
\bottomrule
\end{tabular}

\textbf{观察}：原 Neumann 级数 \emph{不收敛}——\(|a_n|\) 长期保持 \(\mathcal{O}(1)\)，
具体地 \(|a_{10}| \approx 1.55, |a_{11}| \approx 1.43\)，丝毫看不出几何衰减；
\(\rho(\Kker) \gtrsim 1\)。这也是为什么必须 Padé。

\medskip
\textbf{Padé 收敛序列 \(\Pade{N}{N}(z=1)\)}（由 \(\{a_0, \ldots, a_{2N}\}\) 算出，
计算\emph{不依赖}级数收敛，纯代数）：

\medskip
\begin{tabular}{rll}
\toprule
\(N\) & \(\Pade{N}{N}(1)\) & \(\Delta_N := |\Pade{N}{N} - \Pade{N-1}{N-1}|\) \\
\midrule
0 & \(-0.6268 + 0.0000\,i\) & --- \\
1 & \(+0.3115 - 4.5631\,i\) & \(4.66\) \\
2 & \(-0.0681 - 3.2719\,i\) & \(1.36\) \\
3 & \(-0.1209 - 3.1934\,i\) & \(9.4\times 10^{-2}\) \\
4 & \(-0.1217 - 3.1965\,i\) & \(3.2\times 10^{-3}\) \\
5 & \(-0.1217 - 3.1964\,i\) & \(\sim 10^{-5}\) \\
\(\ge 6\) & 同 \(N=5\) 至 6 位有效 & \(< 10^{-7}\)（criterion 3 触发） \\
\bottomrule
\end{tabular}

\textbf{Solver log 中的对应输出}（\texttt{solver\_run.log:495--497}）：
\begin{verbatim}
- Convergence reached for all on-shell elements!
- Extracting on-shell U-matrix elements
  - U-matrix element for alpha'=4, alpha=4, q=15:
       -1.2167241819e-01 + -3.1964688612e+00i
\end{verbatim}
与上表 \(N\ge 5\) 行的 6 位有效完全吻合。

\medskip
\textbf{dense vs.\ Padé 对照}（toy \(\Np=\Nq=10\)，\(J^\pi=1/2^+\) 单通道）：
两条路径解出的 \(\Uampl{\PWchanp=\PWchan=4}{}(q_{\text{shell}})\) 在 \texttt{ZGESV} 与 Padé\(N=4\)
两端均给出 \(-0.1217 - 3.197\,i\)，相对差 \(< 5 \times 10^{-6}\)。
\(\Pade{4}{4}\) 已经达到等价 LAPACK 直接求解的精度——这正是 Padé
作为 dense 替代品的合法性证据。

\medskip
\textbf{时间对比（toy \(D=400\)）}：
\begin{itemize}
  \item dense (\texttt{ZGESV})：\(\sim 2.7\) 秒（含矩阵装配 \(\sim 1.8\) s）；
  \item Padé（\(N=5\) 即收敛）：\(\sim 4.1\) 秒。
\end{itemize}
toy 规模下两者相当；但当 \(D\) 增长到 \(2 \times 10^4\)：
\begin{itemize}
  \item dense：\(\mathcal{O}(D^3) \sim 8 \times 10^{12}\) ops \(\to\) 数小时
        + \(D^2\) 个复数 \(\sim\!6\) GiB 内存（仍可行但已经吃力）；
  \item Padé：\(\mathcal{O}(D^2 N_{\max}) \sim 5.6 \times 10^9\) ops \(\to\) 分钟级
        + \(D \cdot N_{\text{shell}}\) 行 \(\sim\!4\) GiB 内存。
\end{itemize}
当 \(D \to 6 \times 10^4\)（生产规模 \(\Nalpha=64, \Np=\Nq=30\)）：
\begin{itemize}
  \item dense：内存需求 \(\sim 27\) GiB \(\Rightarrow\) 不可行；
  \item Padé：仍可行，约 1.5 小时（与 \cref{ch:09_permutation_p123} 的 \(\Pker\) 缓存共用）。
\end{itemize}
这正是 \cref{sec:ch11-motivation} 中"\(D \lesssim 5000\) 用 dense
否则只能用 Padé"的经验阈值由来。
\end{numericalbox}

% =========================================
\section{接口}
\label{sec:ch11-interfaces}

\paragraph{上游依赖。}
本章求解器站在前面 6 章的肩膀上：
\begin{itemize}
  \item \cref{ch:08_potential_matrix} 提供 \(\Vker\)：\texttt{V\_WP\_unco\_array}, \texttt{V\_WP\_coup\_array}（按 \(\rho\) 索引去重的稀疏 LS 块）；
  \item \cref{ch:09_permutation_p123} 提供 \(\Pker\)：CSC 格式的 \texttt{P123\_sparse\_*\_array}，密度 \(\sim 0.1\%\)；
  \item \cref{ch:10_resolvent} 提供 \(\Gker\)：\texttt{G\_array}，每个 \(q\)-shell 一份 \(D\)-长的对角；
  \item \cref{ch:06_wp_swp_states} 提供 \(\Cker\)：\texttt{swp\_states.C\_SWP\_unco/coup\_array}（FWP→SWP 基变换）；
  \item \cref{ch:07_pw_symm_states} 提供 \texttt{pw\_states} 与 on-shell 通道索引（\texttt{chn\_os\_indexing}）；
  \item \cref{ch:15_3nf_physics} 通过 \texttt{tnf} 指针注入 3NF 贡献——
        默认 \texttt{nullptr} 即 2NF-only，开启后通过 \texttt{calculate\_CPVC\_col} 内部分支累加 \(W^{(1)}\)。
\end{itemize}

\paragraph{下游消费者。}
\(\Uker\) 出口包含两个数组：
\begin{itemize}
  \item \texttt{U\_array}：弹性振幅 \(\Uampl{\PWchanp}{\PWchan}(q_{\text{shell}})\)，
        总长 \(N_{d-\text{row}} \cdot N_{d-\text{col}} \cdot N_{q-\text{com}}\)。
        交给 \cref{ch:12_observables_basics} 计算 \(\dsigma/\dOmega\)、
        \cref{ch:13_polarization_observables} 计算 \(\iTeleven, \Ttwentyzero, \Ttwentyone, \Ttwentytwo\)、
        \cref{ch:14_miller_gate1} 用作相移提取的输入；
  \item \texttt{U\_BU\_array}：破裂振幅（仅在 \texttt{include\_breakup\_channels=true} 时启用）；
        本仓库 190 MeV 工作流默认\textbf{关闭}破裂通道（弹性观测量足够）；
        关闭破裂通道后 \texttt{a\_BU\_coeff\_array=nullptr}，省 \(\sim 30\%\) GEMM 工作量。
\end{itemize}

\paragraph{run\_parameters 中的相关开关。}
\begin{itemize}
  \item \texttt{solve\_dense}：\texttt{true} 切到 dense oracle 路径（\cref{lst:ch11-dense-solver}）；
        默认 \texttt{false}；
  \item \texttt{include\_breakup\_channels}：是否解 BU 振幅；
  \item \texttt{w1\_scale}：3NF 外层缩放因子（\texttt{tnf\_kernel\_context::w1\_scale}），默认 \(1.0\)。
\end{itemize}

% =========================================
\section{常见陷阱}
\label{sec:ch11-pitfalls}

\begin{pitfallbox}
\textbf{陷阱 1：Padé 杂极点（spurious poles）。}
对角 \(\Pade{N}{N}\) 是\textbf{有理}函数 \(P/Q\)，其极点是 \(Q\) 的零点；
若 \(Q\) 的某个零点意外靠近 \(z=1\) 求值点，会导致 \(\Pade{N}{N}(1)\) 出现
\(\sim 10^{6}\) 数量级的伪振幅。这种"杂极点"\textbf{并非物理共振}——
它来自级数的舍入误差通过行列式被放大。

\textbf{识别}：在 \texttt{neumann\_terms\_*.txt} 中
\(|a_n|\) 长期 \(\mathcal{O}(1)\)，但某一阶 \(\Pade{N}{N}(1)\) 突然冲到
\(\mathcal{O}(10^{6})\)，下一阶又跳回 \(\mathcal{O}(1)\)。

\textbf{应对}：本仓库的 \texttt{idx\_best\_PA} 启发式（\cref{lst:ch11-pade-convergence}）
通过\emph{滑动最小差分}选最佳阶，能把杂极点阶剔出去——
但仅当杂极点出现在\emph{不连续}的 \(N\) 时；连续两阶都被杂极点污染就会失败。
此时唯一的正解是\textbf{切换势模型或换 q-shell}（用相邻
on-shell q-bin 重做）。

\medskip
\textbf{陷阱 2：dense\_solver 的内存爆炸。}
设置 \texttt{solve\_dense=true} 时，\texttt{L\_array} 与
\texttt{R\_array} 各占 \(D^2\) 复数。生产规模 \(D \sim 6 \times 10^4\)
单块 \(\sim\!54\) GiB——直接 OOM。运行 dense 前\emph{必须}降到
\(\Np=\Nq \le 12\)，并用 \texttt{two\_J\_3N\_max=1}, \texttt{J\_2N\_max=1}
压缩 \(\Nalpha\)。否则节点会被 Linux OOM killer 强杀，
留下\emph{无任何输出}的失败工况，初学者经常误以为是程序 crash。

\medskip
\textbf{陷阱 3：\(\eps\) 处方对收敛的影响。}
\cref{ch:10_resolvent} 中 \(G_0(z)\) 的 \(\eps\) 处方
（\(z = E + i\eps\)）会渗透到 \(\Gker\) 的虚部。若 \(\eps\) 取得
过小（\(< 10^{-5}\) MeV），近 on-shell 的 \(G\) 元素会出现
\(\mathcal{O}(10^{5})\) 的尖峰，导致 \(\Aker \cdot \Gker\) 的 GEMM 中
舍入误差被放大，进而 Neumann 系数前几项就失精度。
本仓库默认 \(\eps = 10^{-3}\sim 10^{-2}\) MeV，
对 \(\NqWP \ge 20\) 的 SWP 离散化是安全的；若用户手工调小
应同时检查 Neumann 系数的相对噪声。

\medskip
\textbf{陷阱 4：on-shell vs.\ off-shell 提取错误。}
\cref{eq:ch11-on-shell-subset} 要求把 \(p\)-bin 索引固定为 \texttt{elastic\_bound\_packet\_index}
（即 SWP 谱排序后的束缚态索引，\textbf{不是 0}除非约定如此）。
SWP 构造时若 \(\Vker\) 的本征值排序变化（譬如换势从 N2LOopt 到 Idaho\_N3LO），
deuteron 束缚态的 SWP 索引可能从 0 变到其他值。一旦提错就会得到
"看起来合理但物理上错误"的 \(U\)：截面在 \(\theta_{\text{lab}}\)
分布的\emph{形状}对，但\emph{绝对值}差 \(\mathcal{O}(10)\)。
\textbf{诊断方法}：检查 \texttt{solver\_validation\_*.json} 中
\texttt{deuteron\_idx\_array} 与 \texttt{q\_com\_idx\_array}
是否与运行 log 中的"Found bound state with energy ..."一致。

\medskip
\textbf{陷阱 5：commit \texttt{4d120d5} 教训——dense 路径只该用于 debug。}
当年 Sean 引入 \texttt{solve\_dense=true} 时，附文档明确写"slow and costly
but useful for debugging"。然而后续有用户尝试用 dense 跑\(\Tlab\) sweep，
结果运行时间 \(\mathcal{O}(\text{天})\) 量级。\textbf{原则}：
任何\(\Tlab \ge 5\) 个能量点 + \(\Nalpha \ge 30\) 的扫描必须用 Padé；
dense 仅在以下三个场景使用：
(1) 单 \(\Tlab\) 单 \(J^\pi\) 的 oracle 验证；
(2) \(\Aker\) / \(\Lker\) 矩阵的频谱诊断；
(3) commit-bisect 时定位回归 bug。

\medskip
\textbf{陷阱 6：复矩阵实部/虚部拆分对未来 3NF/光学势的限制。}
\cref{eq:ch11-complex-as-real} 的两次 \texttt{dgemm} 优化\textbf{要求 \(\Aker\) 实}。
当前 N2LOopt 等 NN 势都满足；但若引入光学势或带吸收的耦合道势，\(\Aker\)
变为复，需把两次 \texttt{dgemm} 换回一次 \texttt{zgemm}（4 条调用、
若干 BLAS 别名）。在 \cref{lst:ch11-neumann-gemm} 中至少有
\textbf{两处}必改：(a) re/im 数组改为 \texttt{cdouble*}；
(b) GEMM 调用改 \texttt{cblas\_zgemm}（注意 \texttt{alpha/beta} 也要复化）。
不做修改而强行跑光学势会产出"看起来正常但虚部全为 0"的错误结果——
更糟的是它\emph{不会 crash}，只会让虚部默默丢失。
\end{pitfallbox}

% =========================================
\section{小结}
\label{sec:ch11-summary}

本章把 AGS 方程的"形式逆"\((\Iker - \Kker)^{-1}\Aker\) 落到了一台具体可执行的求解器上。
两条路径：
\begin{itemize}
  \item \textbf{Padé 路径}（生产）——把 \(\Uampl{\PWchanp}{\PWchan}\)
        变成一个标量级数 \(\{a_n\}\)，每生成 2 项就升一阶 \(\Pade{N}{N}\)，
        通过 4 个收敛判据自动选最佳阶。原级数发散在物理上是常态，
        Padé 把它解析延拓到真实振幅。
  \item \textbf{Dense 路径}（debug）——直接 \texttt{ZGESV} 解 \((\Iker - \Gker\Aker)\Uker = \Aker\)；
        只能跑 toy 规模，但作为 oracle 不可替代。
\end{itemize}
两条路径共用 \(\Aker = \Cker^\top\Pker\Vker\Cker\) 的列计算（\texttt{calculate\_CPVC\_col}）；
分歧在"如何把 \(\Aker\) 拼成 \(\Uker\)"。

读完本章，你应能：
\begin{enumerate}
  \item 看到一份 \texttt{neumann\_terms\_*.txt} 立即判断该通道是否需要 Padé（看 \(|a_n|\) 是否衰减）；
  \item 给定前 \(2N+1\) 个 \(a_n\)，用 \cref{eq:ch11-pade-det} 手算 \(\Pade{N}{N}(1)\)；
  \item 在 solver log 里读懂 \texttt{Working on Pade approximant P[N,M]}、
        \texttt{Calculating Pade approximants PA[N,M]}、
        \texttt{Convergence reached for all on-shell elements} 三类输出；
  \item 选 dense 还是 Padé，知道每条路径的内存与时间复杂度；
  \item 当一通道的 \(\Pade{N}{N}\) 振荡时识别为杂极点而非物理共振；
  \item 添加新势/新通道时，知道哪些 \cref{lst:ch11-neumann-gemm} 中
        实/复假设需要复审。
\end{enumerate}

下一章（\cref{ch:12_observables_basics}）将把 \texttt{U\_array} 推到
\(\dsigma/\dOmega\) 与极化观测量；本章产出的 \(\Uampl{\PWchanp}{\PWchan}(q)\)
是这一系列下游计算的唯一物理输入。
